\vspace*{-1.5cm}
\section*{Resumen}
\label{sec:resumen}

Este Trabajo Fin de Máster presenta el diseño, implementación y validación de un laboratorio de Red Team orientado a la formación y experimentación en ciberseguridad ofensiva. El entorno reproduce una infraestructura corporativa híbrida compuesta por un servidor Ubuntu con una aplicación web vulnerable y un dominio Active Directory basado en Windows Server 2012 y Windows 7. Sobre esta infraestructura se implementa una cadena completa de ataque que abarca las fases de reconocimiento, explotación web, escalada de privilegios, obtención de credenciales, movimiento lateral y compromiso del dominio.

Durante el desarrollo del proyecto se diseñaron escenarios vulnerables basados en técnicas habituales utilizadas en ejercicios de Red Team, incorporando vulnerabilidades representativas del OWASP Top 10 y técnicas documentadas en el marco MITRE ATT\&CK. Asimismo, se desarrolló un conjunto de scripts y procedimientos que permiten automatizar el despliegue del laboratorio y garantizar su reproducibilidad.

Como resultado, se obtuvo un entorno funcional, modular y reproducible que puede ejecutarse sobre equipos con recursos hardware moderados gracias a un proceso de optimización de memoria y almacenamiento. El laboratorio constituye una plataforma práctica para el aprendizaje de técnicas ofensivas, la realización de pruebas de penetración y el entrenamiento en escenarios similares a los encontrados en infraestructuras corporativas reales.


\vspace{-1.5mm}
\vfill
\section*{Palabras Clave}
\label{sec:palabra_clave}

Red Team, laboratorio de ciberseguridad, Active Directory, pruebas de penetración, OWASP Top 10, MITRE ATT\&CK, escalada de privilegios, movimiento lateral, seguridad ofensiva, virtualización.
\vfill
\vfill
\vfill
\vfill


\cleartooddpage
\clearpage

\vspace*{-1.5cm}
\section*{Abstract}
\label{sec:abstract}
This Master's Thesis presents the design, implementation, and validation of a Red Team laboratory intended for cybersecurity training and offensive security experimentation. The environment reproduces a hybrid corporate infrastructure consisting of an Ubuntu server hosting a vulnerable web application and an Active Directory domain based on Windows Server 2012 and Windows 7. A complete attack chain is implemented, covering the stages of reconnaissance, web exploitation, privilege escalation, credential harvesting, lateral movement, and domain compromise.

Throughout the project, vulnerable scenarios inspired by common Red Team techniques were developed, incorporating representative vulnerabilities from the OWASP Top 10 and attack techniques documented in the MITRE ATT\&CK framework. In addition, deployment scripts and technical documentation were created to automate the installation process and ensure the reproducibility of the environment.

The resulting laboratory is a functional, modular, and reproducible platform that can run on commodity hardware thanks to memory and storage optimization. It provides a practical environment for learning offensive security techniques, performing penetration testing exercises, and training in scenarios that closely resemble real-world corporate infrastructures.




\vfill
\vfill
\section*{Keywords}
\label{sec:keywords}
Red Team, cybersecurity laboratory, Active Directory, penetration testing, OWASP Top 10, MITRE ATT\&CK, privilege escalation, lateral movement, offensive security, virtualization.
\vfill
\vfill
\vfill
\vfill